\par\Needspace{6\baselineskip}\medskip
\SourceLeaf{099}\hypertarget{native-099-article075}{}\pdfbookmark[2]{75.}{article-075}
\begin{center}75.\end{center}
Alia solutio problematis praec. superstruitur observationi in art. 64 III traditae, quod locus heliocentricus terrae, geocentricus corporis coelestis eiusdemque locus heliocentricus in uno eodemque circulo maximo sphaerae sunt siti. Sint in fig. 3 illi loci resp. $T$, $G$, $H$; porro $\AscendingNode$ locus nodi ascendentis; $\AscendingNode T$, $\AscendingNode H$ partes eclipticae et orbitae, $GP$ perpendiculum ad eclipticam ex $G$ demissum, quod igitur erit $=b$. Hinc et ex arcu $PT=L-l$ determinabitur angulus $T$ atque arcus $TG$. Dein in triangulo sphaerico $\AscendingNode HT$ data sunt angulus $\AscendingNode=i$, angulus $T$ latusque $\AscendingNode T=L-\AscendingNode$, unde eruentur duo reliqua latera $\AscendingNode H=u$ atque $TH$. Tandem erit $HG=TG-TH$ atque $r=\frac{R\sin TG}{\sin HG}$, $\Delta=\frac{R\sin TH}{\sin HG}\mathpunct{\cdot}$
